%! Comparing Studies and Choosing a Method — Unit 1, Lesson 1.7 — Group Activity (Answer Key)
\documentclass[10pt]{article}
\usepackage{saar-article}
\usepackage{saar-key}   % pulls in -boxes; do NOT also load -boxes
\newcommand{\TallMath}[1]{$\displaystyle #1\rule[-1.4em]{0pt}{3.2em}$}

\begin{document}

\pageheader{Unit 1, Lesson 1.7}{Group Activity --- Answer Key}
% NAMESTRIP: no \namedateperiod here — mirrors the blank. Teacher-only prose goes
% in the lesson plan as \begin{teachernote}[Group Activity], never in a key.

\vspace{0.15in}

\begin{headlinebox}{frost}
\small
\textbf{Three Studies on One Desk --- Harbor Point Community Pool.} The director
has run two studies on her $1{,}500$ members already. \textbf{In Lesson 1.5 she
surveyed:} a random sample of $200$ members, $70\%$ of the morning swimmers
against $45\%$ of the evening swimmers slept at least $7$ hours --- a $25$-point
gap, but the members had picked their own swim times. \textbf{In Lesson 1.6 she
assigned:} $90$ volunteers split $45$/$45$ by a generator, $60\%$ against
$40\%$ --- a $20$-point gap. \textbf{Now the board asks a new question:} does a
weekly water-aerobics class reduce reported joint pain? She takes
\textbf{$120$ volunteers}, lets a generator send \textbf{$60$ to the class} and
\textbf{$60$ to no class} for six weeks, and keeps everything else the same. At
the end, \textbf{$42$} of the class group and \textbf{$30$} of the no-class
group reported less joint pain.
\end{headlinebox}

\vspace{0.03in}

\boxguard[30]
\begin{tcolorbox}[colback=white, colframe=black!40, breakable,
    title=\textbf{Tier R --- Remediate},
    fonttitle=\bfseries, arc=2mm, left=3mm, right=3mm, top=1mm, bottom=1mm]
\small
\begin{enumerate}[leftmargin=1.6em, itemsep=2pt, topsep=3pt]

  \item Label the director's two earlier studies.

        The Lesson 1.5 swim-time survey: \ans{observational study}

        The Lesson 1.6 swim-time study: \ans{controlled experiment}

  \item Who decided which group each person was in?

        \textbf{(a)} In the survey: \ans{the members --- they chose their times}

        \textbf{(b)} In the assigned study: \ans{chance --- a random generator}

  \item In the new water-aerobics experiment, $42$ of the $60$ in the class
        reported less joint pain. What percent is that?

        \begin{work}
          \dfrac{42}{60} &= 0.70 \\
                         &= 70\%
        \end{work}

  \item $30$ of the $60$ with no class reported less joint pain. What percent is
        that?

        \begin{work}
          \dfrac{30}{60} &= 0.50 \\
                         &= 50\%
        \end{work}

  \item Which group did better, and by how many percentage points?
        \ansline{The class group --- $70\% - 50\% = 20$ percentage points higher.}
\end{enumerate}
\end{tcolorbox}

\vspace{0.03in}

\boxguard
\begin{tcolorbox}[colback=white, colframe=black!40, breakable,
    title=\textbf{Tier A --- Approaching Proficiency},
    fonttitle=\bfseries, arc=2mm, left=3mm, right=3mm, top=1.5mm, bottom=1.5mm]
\small
\begin{enumerate}[leftmargin=1.6em, itemsep=3pt, topsep=3pt]

  \item Put the water-aerobics volunteers back together.

        \textbf{(a)} How many of the $120$ reported less joint pain, and what
        percent is that?

        \begin{work}
          42 + 30 &= 72 \text{ volunteers} \\
          \dfrac{72}{120} &= 0.60 = 60\%
        \end{work}

        \textbf{(b)} About how many of all $1{,}500$ members does that predict?

        \begin{work}
          0.60 \times 1500 &= 900 \text{ members}
        \end{work}

  \tcbbreak   % keep item 2's prompt with its table — mirrored in activity_key
  \item Fill in the plan the director followed, one stage at a time.

        \begin{center}
        \renewcommand{\arraystretch}{1.4}
        \begin{tabularx}{0.95\linewidth}{@{} p{2.4cm} X @{}}
        \toprule
        \textbf{Stage} & \textbf{The water-aerobics experiment} \\
        \midrule
        Formulate & Does a weekly water-aerobics class reduce reported joint
          pain among Harbor Point members? \\
        Collect & \ans{Take $120$ volunteers; a generator assigns $60$ to the
          weekly class and $60$ to no class for six weeks.} \\
        Represent & \ans{A two-row table --- each group's size, how many
          reported less joint pain, and each group's percent.} \\
        Analyze \& report & \ans{$70\%$ against $50\%$, a $20$-point gap; say
          what the class did for these members.} \\
        \bottomrule
        \end{tabularx}
        \end{center}

  \item Write the strongest sentence each study earns.

        \textbf{(a)} The Lesson 1.5 swim-time survey:
        \ansline{Swimming in the morning is associated with more sleep --- an association only.}

        \textbf{(b)} The water-aerobics experiment:
        \ansline{The weekly water-aerobics class caused less reported joint pain.}

  \item Some volunteers quit the class after two weeks and the director wants to
        know \emph{why}. Name the collection method you would use, and say why a
        written survey is a poor fit for this particular question.
        \ansline{Interviews. ``Why'' needs depth and follow-up questions, and only a few people quit.}
        \ansline{A fixed survey gives every quitter the same boxes and cannot ask the next question.}
\end{enumerate}
\end{tcolorbox}

\vspace{0.03in}

\boxguard
\begin{tcolorbox}[colback=white, colframe=black!40, breakable,
    title=\textbf{Tier E --- Extension},
    fonttitle=\bfseries, arc=2mm, left=3mm, right=3mm, top=1.5mm, bottom=1.5mm]
\small
\begin{enumerate}[leftmargin=1.6em, itemsep=4pt, topsep=3pt]

  \item The board asks a fourth question: \emph{do members who already have
        arthritis sleep worse than members who do not?}

        \textbf{(a)} Could the director assign the groups? \ans{No}

        \textbf{(b)} Say why or why not, and name the study she must run
        instead.
        \ansline{Arthritis is not a treatment you can hand out --- so it must be observational.}

  \item Suppose that study finds arthritis members sleep worse. Write the
        strongest claim she may publish, and name one other variable that could
        produce that difference on its own.
        \ansline{Having arthritis is associated with less sleep at Harbor Point --- not ``causes.''}
        \ansline{Age: older members are likelier to have arthritis and likelier to sleep poorly.}

  \item Name the best collection method for each job.

        \textbf{(a)} Count how many of all $1{,}500$ members would sign up for a
        class: \ans{a survey}

        \textbf{(b)} Understand in depth why three specific members quit:
        \ans{interviews}

        \textbf{(c)} Record whether members actually use the handrail when
        entering the pool: \ans{observation}
\end{enumerate}
\end{tcolorbox}

\end{document}
